% =====================================================================
% Real-world (HSLU) generalisation subsection for Chapter 5.
%
% Insert inside \section{Technical Evaluation of Data Fusion}
% (\label{sec:technical-evaluation}), immediately AFTER the existing
% \subsection{Discussion of Technical Findings} and BEFORE
% \section{Usability Evaluation of the Adaptive Interface}.
%
% Numbers are from cmd/fusion-bench: 62 HSLU households, 60 s dataset-time
% grid, 14-day window per household, both predictors run as-is (no retuning).
% Reproduce with:  cd backend && make bench-census ; make bench-all
% =====================================================================

\subsection{Generalisation to Real-World Data}
\label{sec:hslu-evaluation}

The evaluation in Section~\ref{sec:technical-evaluation} was conducted under
controlled conditions: the simulator provided a known expected label for every
inference cycle, which made it possible to measure accuracy directly. To examine
whether the two predictors behave sensibly when exposed to genuine, uncontrolled
sensor streams, a second evaluation was carried out on a public real-world
dataset.

\subsubsection{Dataset and Setup}

The HSLU smart-home dataset\footnote{\url{https://doi.org/10.5281/zenodo.14243471}}
records ambient conditions in 62 households across four European countries,
predominantly inhabited by older adults, over a period of 359 days. It contains
periodic measurements of temperature, humidity and ambient light sampled every
two to five minutes, together with event-based passive-infrared (PIR) motion and
door-contact signals. The raw exports were preprocessed into one chronologically
ordered timeline per household, retaining the sensor types used by the fusion
models (temperature, humidity, ambient light and motion).

Two characteristics of the dataset shaped the evaluation. First, it contains
\emph{no} activity or context annotations. The expected-label accuracy reported
in Table~\ref{tab:fusion-results} therefore cannot be computed on it: there is no
ground truth against which a prediction could be marked correct. Second, the PIR
stream records only activations and never deactivations, so presence cannot be
read from an explicit ``off'' signal. Accordingly, motion presence was modelled
by decaying from the most recent activation in \emph{dataset} time, which is the
same recency mechanism the predictors already use, evaluated against the
timestamp of each decision rather than wall-clock time.

Given the absence of ground truth, the evaluation is deliberately
\emph{label-free}. Rather than measuring correctness, it characterises how the
two predictors \emph{behave} when transferred, without any retuning, from the
synthetic calibration of Section~\ref{sec:technical-evaluation} to real data. An
offline benchmark replayed each household's timeline on a fixed 60-second grid in
dataset time; at every grid point the latest reading per sensor formed a sensor
window that was passed to both predictors. A representative 14-day window per
household was used. The following properties were recorded:

\begin{itemize}
    \item \textbf{Inter-model agreement} -- the percentage of decisions for which
    the rule-based and fuzzy predictors returned the same label.
    \item \textbf{Context distribution} -- the share of decisions assigned to each
    context label by each predictor.
    \item \textbf{Label stability} -- the number of label transitions per hour and
    the number of single-tick flips (a label change immediately reverted on the
    next decision), measured on the raw output of each predictor.
    \item \textbf{Confidence} and \textbf{inference latency} per predictor.
    \item \textbf{Smoothing effect} -- the reduction in transitions after applying
    the same temporal smoother used in production ($N = 2$ consecutive
    confirmations) to each raw label stream.
\end{itemize}

\subsubsection{Sensor Coverage}

The predictors require particular sensor channels in particular rooms: living-room
temperature, humidity, light and motion drive the comfort, alert, reading and
television contexts, while the cooking and sleeping contexts additionally require
kitchen and bedroom light and motion. A census of all 62 households
(Table~\ref{tab:hslu-coverage}) shows that real deployments are frequently
under-instrumented: many homes report ambient light in only one room, and three
homes record neither temperature nor humidity. The benchmark was therefore run on
the whole population and on two coverage-filtered subsets, so that the behavioural
comparison is not distorted by households in which entire context classes can
never fire.

\begin{table}[htbp]
\captionsetup{justification=centering}
\caption{Sensor coverage across the 62 HSLU households. ``Usable'' and ``full''
define the coverage-filtered subsets used for the behavioural comparison.}
\label{tab:hslu-coverage}
\centering
\small
\renewcommand{\arraystretch}{1.2}
\begin{tabular}{@{}p{0.18\textwidth} p{0.56\textwidth} p{0.16\textwidth}@{}}
\hline
\textbf{Tier} & \textbf{Criteria} & \textbf{Households} \\
\hline
All recorded & at least one supported sensor channel & 62 \\
Usable & temperature and humidity present, plus living-room light and motion & 19 \\
Full & usable, plus kitchen and bedroom light and motion & 10 \\
\hline
\end{tabular}
\end{table}

\subsubsection{Results}

Across all coverage levels the two predictors agreed on roughly four decisions in
five, and both produced very similar label-transition rates
(Table~\ref{tab:hslu-subsets}). This is a higher agreement than the 71.5\%
observed on the simulator, and the reason is visible in the context distribution:
real elderly-resident homes are dominated by long unoccupied periods, so the
absence context accounts for the majority of all decisions, and absence is a state
on which the two models trivially concur.

\begin{table}[htbp]
\captionsetup{justification=centering}
\caption{Behaviour of the two predictors across the full population and the two
coverage-filtered subsets (60 s grid, 14-day window per household).}
\label{tab:hslu-subsets}
\centering
\small
\renewcommand{\arraystretch}{1.2}
\begin{tabular}{@{}p{0.26\textwidth} p{0.13\textwidth} p{0.16\textwidth} p{0.18\textwidth} p{0.18\textwidth}@{}}
\hline
\textbf{Subset} & \textbf{Homes} & \textbf{Agreement} & \textbf{Rule trans./h} & \textbf{Fuzzy trans./h} \\
\hline
All recorded & 62 & 80.5\% & 1.65 & 2.04 \\
Usable       & 19 & 83.7\% & 1.96 & 1.95 \\
Full         & 10 & 81.4\% & 1.83 & 1.83 \\
\hline
\end{tabular}
\end{table}

Table~\ref{tab:hslu-behaviour} reports the detailed comparison on the
full-coverage subset, in which every context the models can emit is structurally
reachable. Both predictors completed inference in under two microseconds and
remained far below the one-second bound required for real-time operation. The
fuzzy predictor was again marginally slower and reported markedly higher mean
confidence. The two models produced an essentially identical number of label
transitions per hour, but the fuzzy predictor generated almost twice as many
single-tick flips on the raw stream; the production temporal smoother removed a
comparable fraction of transitions for both.

\begin{table}[htbp]
\captionsetup{justification=centering}
\caption{Behavioural comparison on the full-coverage subset (10 households).
Latency and confidence are means over all decisions; flips are totals.}
\label{tab:hslu-behaviour}
\centering
\small
\renewcommand{\arraystretch}{1.2}
\begin{tabular}{@{}p{0.46\textwidth} p{0.22\textwidth} p{0.22\textwidth}@{}}
\hline
\textbf{Metric} & \textbf{Rule-Based} & \textbf{Fuzzy Logic} \\
\hline
Mean inference latency & 1.77~$\mu$s & 1.89~$\mu$s \\
Inference latency (95th percentile) & 2.38~$\mu$s & 2.67~$\mu$s \\
Mean confidence & 0.864 & 0.955 \\
Label transitions per hour & 1.83 & 1.83 \\
Single-tick flips (total) & 148 & 279 \\
Flicker reduction after smoothing & 11.8\% & 12.8\% \\
\hline
\end{tabular}
\end{table}

The clearest differences appear in the context distribution
(Table~\ref{tab:hslu-distribution}). Both models assign the majority of decisions
to \texttt{no\_one\_home}. They diverge most on temperature interpretation: the
fuzzy predictor labels \texttt{alert\_too\_cold} in roughly 9\% of decisions,
whereas the rule-based predictor almost never does. They also diverge on activity
contexts, where the rule-based predictor fires \texttt{cooking\_kitchen} and
\texttt{watching\_tv} roughly twice as often as the fuzzy predictor.

\begin{table}[htbp]
\captionsetup{justification=centering}
\caption{Context-label distribution on the full-coverage subset (10 households),
as a percentage of all decisions per predictor.}
\label{tab:hslu-distribution}
\centering
\small
\renewcommand{\arraystretch}{1.2}
\begin{tabular}{@{}p{0.40\textwidth} p{0.22\textwidth} p{0.22\textwidth}@{}}
\hline
\textbf{Context} & \textbf{Rule-Based} & \textbf{Fuzzy Logic} \\
\hline
\texttt{no\_one\_home}    & 62.1\% & 63.8\% \\
\texttt{comfortable}      & 16.7\% & 10.4\% \\
\texttt{cooking\_kitchen} & 11.8\% & \phantom{0}6.7\% \\
\texttt{alert\_too\_hot}  & \phantom{0}5.1\% & \phantom{0}8.1\% \\
\texttt{alert\_too\_cold} & \phantom{0}0.0\% & \phantom{0}8.8\% \\
\texttt{watching\_tv}     & \phantom{0}4.2\% & \phantom{0}1.7\% \\
\texttt{reading}          & \phantom{0}0.1\% & \phantom{0}0.2\% \\
\texttt{sleeping}         & \phantom{0}0.0\% & \phantom{0}0.0\% \\
\hline
\end{tabular}
\end{table}

\subsubsection{Discussion of Real-World Findings}

The real-world evaluation complements, rather than overturns, the simulator
result, and it exposes three effects that a controlled simulator cannot.

First, the high inter-model agreement is a property of the deployment rather than
of the algorithms: it arises because the households are unoccupied for most of the
recording, and absence is unambiguous. The informative disagreement is therefore
concentrated in the minority of occupied or boundary states -- precisely the
regime in which the choice of fusion strategy matters.

Second, the divergence on \texttt{alert\_too\_cold} is a transfer artefact of
calibration. The fuzzy ``cold'' membership assigns a non-zero degree from
approximately $19\,^{\circ}\text{C}$ downward, so the many real homes that sit in
the $17$--$19\,^{\circ}\text{C}$ range are partially classified as cold, whereas
the rule-based predictor's hard $17\,^{\circ}\text{C}$ threshold is almost never
crossed.
The membership functions and thresholds were tuned to the synthetic value ranges
of the simulator; when applied unchanged to real distributions they remain
internally consistent but no longer reflect the same lived comfort boundaries.
This indicates that, for deployment, the linguistic terms should be recalibrated
to the temperature and illuminance distributions actually observed in the target
homes.

Third, the stability advantage of the fuzzy predictor observed near simulator
thresholds does not transfer directly. On real multi-room noise the fuzzy
argument-maximum switches more frequently at the single-sample level, producing
about twice as many single-tick flips as the rule-based predictor, even though the
overall transition rate is comparable. After the production temporal smoother both
predictors reach a similar stability, which confirms that single-sample
robustness in the delivered system is provided primarily by the smoother rather
than by the fusion algorithm alone.

With respect to \textbf{RQ1}, the real-world evaluation shows that both fusion
strategies transfer to genuine heterogeneous sensor data and remain comfortably
within real-time limits, and that they agree on clearly defined states while
diverging in ambiguous ones. It also shows that the quality of inference on real
data depends not only on the fusion mechanism but on calibrating that mechanism to
the deployment's sensor coverage and value distributions -- a dependency that the
labelled simulator, by construction, cannot reveal.
